\section{Geometric residuals and prime-threshold components}\label{sec:geometry}

\subsection{Base volume and height}

Let $H_n$ be the $(n-1)$ by $n$ matrix consisting of rows two through $n$ of $A_n$, equivalently of $B_n$. It has full row rank because its minor on columns two through $n$ is unit lower triangular. Proposition~\ref{def:inverse} gives
\[
 H_n\muvec_n=0,\qquad \ker H_n=\spanop\{\muvec_n\}.
\]
The vector of signed maximal minors of $H_n$, normalized to have first coordinate one, is therefore $\muvec_n$. By Cauchy--Binet,
\begin{equation}\label{geo:base}
 \det(H_nH_n^T)=\sum_{j=1}^n\mu(j)^2=Q(n).
\end{equation}
The parallelotope generated by those rows has $(n-1)$-dimensional volume $\sqrt{Q(n)}$. The distance of the remaining row $\one_n^T$ from their span is
\begin{equation}\label{geo:height}
 h_n=\frac{|\one_n^T\muvec_n|}{\norm{\muvec_n}_2}
     =\frac{|M(n)|}{\sqrt{Q(n)}},\qquad
 \min_{z\in\mathbb R^{n-1}}\norm{\one_n-H_n^Tz}_2^2
     =\frac{M(n)^2}{Q(n)}.
\end{equation}
This is an exact height formula, not an estimate for that height.

The symmetric positive semidefinite matrix $G_n^{\mathrm{row}}=H_n^TH_n$ has a useful graph interpretation. On the squarefree coordinates it is the signless Laplacian of the tree: each edge contributes $(e_i+e_j)(e_i+e_j)^T$. Multiplication on both sides by the diagonal matrix with entries $\mu(i)$ on these coordinates changes each edge sum to an edge difference, hence gives the ordinary graph Laplacian. Each nonsquarefree coordinate contributes a one-dimensional identity block.

For $\tau>0$, spectral decomposition yields
\begin{equation}\label{geo:resolvent}
 h_n^2\le\tau\one_n^T(G_n^{\mathrm{row}}+\tau I_n)^{-1}\one_n,
 \qquad
 \lim_{\tau\downarrow0}\tau\one_n^T(G_n^{\mathrm{row}}+\tau I_n)^{-1}\one_n=h_n^2.
\end{equation}
Indeed a spectral coordinate of eigenvalue $\lambda\ge0$ is weighted by $\tau/(\lambda+\tau)$, which equals one on the kernel and tends to zero on each positive eigenvalue. Thus a regularized ellipsoid gives an upper certificate, but merely letting $\tau$ decrease to zero adds no arithmetic information. The weights of $\one_n$ in the eigenspaces matter, not just the eigenvalues.

\subsection{Removing edges by prime size}

Fix a real $y\ge2$. From $H_n$ retain every nonsquarefree diagonal row, and retain a squarefree edge row only if its added prime is greater than $y$. Call the resulting rectangular matrix $H_{n,y}$. Its squarefree components have roots
\[
 \mathcal A(n,y)=\{a\le n:\mu(a)^2=1,\ P^+(a)\le y\}
\]
and vertex sets
\[
 \mathcal C_a(n,y)=\{ad\le n:\mu(d)^2=1,\ P^-(d)>y\}.
\]
To see this, keep the primes at most $y$ in a vertex fixed and delete its larger primes in descending order. The process reaches exactly one root $a$, and reverses along retained edges. The components partition the squarefree vertices.

Set
\begin{equation}\label{geo:F}
 N_a=\Phi_{\sfq}(n/a,y),\quad S_a=\mu(a)R(n/a,y),\quad
 F(n,y)=\sum_{a\in\mathcal A(n,y)}\frac{R(n/a,y)^2}{\Phi_{\sfq}(n/a,y)}.
\end{equation}
Each denominator is a positive component size. Thus $F$ is a sum of squared signed component sums divided by component sizes. It is different from both sums in \eqref{def:energies}.

\begin{proposition}[Projection and variance]\label{geo:projection}
For $n\ge2$ and $y\ge2$,
\begin{align}
 \min_{z\in\mathbb R^{n-|\mathcal A(n,y)|}}\norm{\one_n-H_{n,y}^Tz}_2^2&=F(n,y),\label{geo:residual}\\
 F(n,y)&=\frac{M(n)^2}{Q(n)}+
  \sum_{a\in\mathcal A(n,y)}N_a
       \left(\frac{S_a}{N_a}-\frac{M(n)}{Q(n)}\right)^2.
       \label{geo:variance}
\end{align}
In particular $M(n)^2/Q(n)\le F(n,y)\le Q(n)$. For fixed $n$, $F(n,y)$ is nondecreasing in $y$. For $y\ge\sqrt n$,
\begin{equation}\label{geo:plateaubound}
 0\le Q(n)-F(n,y)\le\frac{4n}{y}.
\end{equation}
\end{proposition}
\begin{proof}
On the component $\mathcal C_a$ the kernel of the retained edge rows is spanned by the alternating vector whose entries are $\mu(ad)=\mu(a)\mu(d)$. Its squared norm is $N_a$, and its inner product with the all-ones vector is $S_a$. Orthogonal projection onto this one-dimensional kernel therefore has squared norm $S_a^2/N_a$. Distinct components have disjoint supports. Nonsquarefree diagonal rows have zero kernel on their coordinates. Summing proves \eqref{geo:residual}.

The partition gives $\sum_aN_a=Q(n)$ and $\sum_aS_a=M(n)$. Expanding the squares in \eqref{geo:variance} proves the equality, including the factor $\mu(a)$ in $S_a$. Omitting that factor would give an incorrect variance identity. The upper bound follows from $|S_a|\le N_a$.

Increasing $y$ deletes edges and splits components. If signed sums and sizes in the pieces are $S_i,N_i$, Cauchy--Schwarz gives $\sum_iS_i^2/N_i\ge(\sum_iS_i)^2/\sum_iN_i$. Thus the energy cannot decrease. If $y\ge\sqrt n$, a rough part is either one or a single prime. A component of size $N$ has alternating sum, up to sign, $2-N$ and energy $(N-2)^2/N=N-4+4/N$. Its deficit from its size is between zero and four. Let $q_1$ be the smallest prime greater than $y$. A nonsingleton root satisfies $a\le n/q_1$, so there are at most $n/q_1$ such components. Summing the deficits proves \eqref{geo:plateaubound}.
\end{proof}

\subsection{A fixed-power profile}

The Dickman function $\rho$ is the continuous function satisfying $\rho(t)=1$ on $0\le t\le1$ and $t\rho'(t)=-\rho(t-1)$ for $t>1$. The Buchstab function $\omega$ is continuous on $[1,\infty)$, equals $1/t$ on $[1,2]$, and satisfies $(t\omega(t))'=\omega(t-1)$ for $t>2$. The symbol $\omega$ in this section denotes this function, not a number of prime factors. Their delay equations determine them successively on unit intervals.

We use the following classical fixed-parameter counting inputs. As $Y\to\infty$, for each fixed $t>0$,
\begin{equation}\label{geo:smoothinput}
 \frac{\Psi(Y^t,Y)}{Y^t}\longrightarrow\rho(t).
\end{equation}
The same limit holds after multiplying $Y^t$ by any fixed positive constant; this follows by monotone bracketing and continuity of $\rho$. Uniformly when $v=\log x/\log y$ belongs to a compact subset of $(1,\infty)$,
\begin{equation}\label{geo:roughinput}
 \Phi(x,y)=\frac{x}{\log y}(\omega(v)+o(1)),\qquad y\to\infty.
\end{equation}
Smooth counting is reviewed in~\cite{hildebrandtenenbaum,mcnew}. For the rough formula on $v\ge2$ see Theorem~B of~\cite{fanpomerance}; on a compact subset of $(1,2]$ the formula follows from $\Phi(x,y)=1+\pi(x)-\pi(y)$ and the prime number theorem. We also use $\Phi(x,y)\ll_V x/\log y$ for $x\ge y$ and $\log x/\log y\le V$, including the endpoint strip. For $1\le v\le2$ this follows from Chebyshev's bound and $1\ll x/\log y$; above two it follows from the rough formula on compact intervals.

\begin{theorem}[Component profile and cutoff]\label{geo:profile}
For each fixed $u\ge1$,
\begin{equation}\label{geo:J}
 \frac{F(n,n^{1/u})}{n}\longrightarrow c_0J(u),\qquad
 J(u)=\rho(u)+\int_1^u\rho(u-v)\frac{\rho'(v)^2}{\omega(v)}\,dv.
\end{equation}
The right derivative of $\rho$ at one may be used at the immaterial endpoint. The function $J$ is positive and nonincreasing, equals one on $[1,2]$, and tends to zero as $u\to\infty$. For $2\le u\le3$,
\begin{equation}\label{geo:explicitJ}
 J(u)=1-4\int_2^u\frac{\log(v-1)}{v(1+\log(v-1))}\,dv,
 \qquad J(2+h)=1-h^2+O(h^3).
\end{equation}
For an arbitrary real sequence $y(n)\ge2$,
\begin{equation}\label{geo:cutoff}
 F(n,y(n))=o(n)\quad\Longleftrightarrow\quad
 \frac{\log y(n)}{\log n}\longrightarrow0.
\end{equation}
The last assertion follows by monotonicity from fixed-$u$ limits; no uniform formula for a growing $u$ is asserted.
\end{theorem}

\begin{proof}
We give the counting and summation details, since the singleton components contribute a term of order $n$ and cannot be dropped.

First define $\Psi_{\sfq}(z,Y)$ to count squarefree $Y$-smooth integers at most $z$. The identity for $\mu^2$ gives
\[
 \Psi_{\sfq}(z,Y)=\sum_{\substack{d\le\sqrt z\\P^+(d)\le Y}}
                         \mu(d)\Psi(z/d^2,Y).
\]
For $z=Y^t$ and each fixed $d$, the normalized summand tends to $\mu(d)\rho(t)/d^2$. The absolute tail beyond $d>D$ is at most $\sum_{d>D}d^{-2}$. Therefore
\begin{equation}\label{geo:sfsmooth}
 \Psi_{\sfq}(Y^t,Y)/Y^t\longrightarrow c_0\rho(t).
\end{equation}
Place mass $1/(a\log Y)$ at $\log a/\log Y$ for each squarefree $Y$-smooth integer $a$, and call this measure $\nu_Y$. Partial summation shows, for fixed $t>0$,
\[
 \nu_Y([0,t])=
 \frac{\Psi_{\sfq}(Y^t,Y)}{Y^t\log Y}
 +\frac1{\log Y}\int_1^{Y^t}\frac{\Psi_{\sfq}(z,Y)}{z^2}\,dz
 \longrightarrow c_0\int_0^t\rho(w)\,dw.
\]
For the integral substitute $z=Y^w$ and use \eqref{geo:sfsmooth} and domination by one. Thus on every compact interval $\nu_Y$ converges weakly to the measure with density $c_0\rho(w)$. The atom at $a=1$ vanishes because its mass is $1/\log Y$.

Second, we derive the signed counterpart of \eqref{geo:roughinput} in the fixed-parameter range:
\begin{equation}\label{geo:signedfixed}
 R(x,y)=\frac{x}{\log y}(\rho'(v)+o(1))
\end{equation}
uniformly on compact subsets of $1<v<\infty$. For $1+\eta\le v\le2$, there are only one and prime rough terms, so $R=1-\pi(x)+\pi(y)$ and the limiting coefficient is $-1/v$. For larger $v$, grouping terms by their least prime factor gives the exact identity
\[
 R(x,y)=R(x,\sqrt x)-\sum_{y<p\le\sqrt x}R(x/p,p).
\]
Proceed by induction over a bounded range of $v$. The inner parameter $w=\log(x/p)/\log p$ is at most $v-1$. In its strip $1\le w<1+\epsilon$ use $|R|\le\Phi\ll x/(p\log p)$. Chebyshev and partial summation, or the prime number theorem, bound the strip by $O(\epsilon x/\log x)+o(x/\log x)$. Off this strip the inductively uniform limit may be summed, since the total of the weights $x/(p\log p)$ is $O_V(x/\log x)$. Prime partial summation gives
\[
 \frac{x}{\log x}\int_1^{v-1}r(w)\,dw
\]
for the sum's leading term, where $r$ is the previously determined coefficient. All primes in these bounded parameter ranges tend to infinity. Since $R(x,\sqrt x)\sim-x/\log x$, the recurrence for the coefficient is
\[
 vr(v)=-1-\int_1^{v-1}r(w)\,dw.
\]
It agrees with $r(v)=\rho'(v)$ by the Dickman equation and the base interval. The recurrence determines each next unit interval, proving \eqref{geo:signedfixed}; first take $x\to\infty$, then $\epsilon\downarrow0$ to remove the strip.

Nonsquarefree rough integers contribute at most $x\sum_{p>y}p^{-2}=O(x/y)$. Hence \eqref{geo:roughinput} also holds for $\Phi_{\sfq}$, with the same positive coefficient. Put $g(v)=\rho'(v)^2/\omega(v)$. We have proved, uniformly on the same compact subsets,
\begin{equation}\label{geo:componentasymptotic}
 \frac{R(x,y)^2}{\Phi_{\sfq}(x,y)}
       =\frac{x}{\log y}(g(v)+o(1)).
\end{equation}

Now fix $u>1$ and put $y=n^{1/u}$. Roots $a>n/y$ have $n/a<y$, so their components are singletons and contribute
\[
 \Psi_{\sfq}(n,y)-\Psi_{\sfq}(n/y,y)=c_0n\rho(u)+o(n).
\]
For the other roots write $v=u-\log a/\log y$. On $v\ge1+\epsilon$, apply \eqref{geo:componentasymptotic} and the weak limit of $\nu_y$ to obtain, after division by $n$,
\[
 \frac1{\log y}\sum_{\substack{a\in\mathcal A(n,y)\\a\le n/y^{1+\epsilon}}}
       \frac{g(u-\log a/\log y)}a
 \longrightarrow c_0\int_{1+\epsilon}^u\rho(u-v)g(v)\,dv.
\]
The summed uniform remainder is $o(1)$ because $(\log y)^{-1}\sum_{a\le n}a^{-1}=O_u(1)$. Endpoint discontinuities in the cutoff cause no problem because the limiting measure has no atoms. In the omitted strip $1\le v<1+\epsilon$, the bound
\[
 R(x,y)^2/\Phi_{\sfq}(x,y)\le\Phi_{\sfq}(x,y)\le\Phi(x,y)\ll x/\log y
\]
gives $O(\epsilon n)+o(n)$ after summing over $y^{u-1-\epsilon}<a\le y^{u-1}$. Choose $0<\epsilon<u-1$; the harmonic sum of this interval is at most $\epsilon\log y+o(1)$. Let $\epsilon\downarrow0$ to prove \eqref{geo:J}. For $u=1$, every component is a singleton and the formula is immediate.

Positivity of $J$ follows from $\rho(u)>0$. Monotonicity follows from Proposition~\ref{geo:projection}, since $n^{1/u}$ decreases with $u$. For completeness the delay equation gives $u\rho(u)=\int_{u-1}^u\rho(t)\,dt$, which prevents a first zero of $\rho$. It also gives monotone decrease and limit zero: a positive limiting value would contradict integrating $-\rho'(u)\ge c/u$. The Buchstab equation gives $\omega(v)\ge1/2$ on $[1,\infty)$ by induction from $v\omega(v)=1+\int_2^v\omega(t-1)\,dt$. Thus $g(v)\le2/v^2$. Dominated convergence in the integral in \eqref{geo:J} proves $J(u)\to0$.

Repeat the component summation with the masses $\Phi_{\sfq}$ instead of squared sums divided by masses. Their total is $Q(n)$, so
\begin{equation}\label{geo:partitionlimit}
 1=\rho(u)+\int_1^u\rho(u-v)\omega(v)\,dv.
\end{equation}
For $1\le v\le2$, $g(v)=\omega(v)=1/v$, giving the plateau. For $2\le u\le3$, only $2\le v\le u$ contributes to the difference between \eqref{geo:partitionlimit} and \eqref{geo:J}, and $\rho(u-v)=1$ there. Directly from the delay equations,
\[
 \rho'(v)=\frac{\log(v-1)-1}{v},\quad
 \omega(v)=\frac{1+\log(v-1)}v,\quad
 \omega(v)-g(v)=\frac{4\log(v-1)}{v(1+\log(v-1))}.
\]
This proves \eqref{geo:explicitJ}; its integrand including the factor four is $2(v-2)+O((v-2)^2)$ at two.

Finally, if $\log y(n)/\log n\to0$, for each fixed $U\ge1$ the cutoff is eventually at most $n^{1/U}$. Monotonicity gives $\limsup F(n,y(n))/n\le c_0J(U)$, which tends to zero as $U\to\infty$. If that logarithmic ratio does not tend to zero, some subsequence satisfies $y(n)\ge n^\delta$ for a fixed $0<\delta<1$. Along it the lower bound tends to $c_0J(1/\delta)>0$. This proves both directions of \eqref{geo:cutoff}.
\end{proof}

\begin{proposition}[Fixed thresholds and elementary matching]\label{geo:limitations}
For each fixed $y\ge2$, $F(n,y)=o(n)$ is equivalent to $M(n)=o(n)$. A matching consisting of retained squarefree tree edges contains at most $\lfloor n/q_1\rfloor$ edges, where $q_1$ is the least prime greater than $y$. Consequently such a matching leaves at least $Q(n)-2\lfloor n/q_1\rfloor$ squarefree vertices uncovered.
\end{proposition}
\begin{proof}
One direction of the equivalence follows from \eqref{geo:variance} and $Q(n)\le n$. Conversely, the finite convolution
\[
 R(x,y)=\sum_{\substack{b\le x\\P^+(b)\le y}}M(x/b)
\]
and $\sum_{P^+(b)\le y}b^{-1}=\prod_{p\le y}(1-1/p)^{-1}<\infty$ imply $R(x,y)=o(x)$ by dominated convergence if $M(x)=o(x)$. For fixed $y$ there are only finitely many squarefree smooth roots. Each denominator is asymptotic to a positive constant times $n/a$, by elementary squarefree density with finitely many excluded primes. Each contribution is therefore $o(n)$ and so is their finite sum. In a matching the smaller endpoints are distinct. Each lies at most $n/q_1$, because its edge adds a prime at least $q_1$. This bounds the number of edges and proves the uncovered-vertex assertion.
\end{proof}

The matching bound limits this particular pairing certificate; it does not lower-bound $F$ itself. More generally a star has the same sparse edge-row form and a single alternating kernel vector, yet its alternating sum can have linear size. Volume, sparsity, and nested tree structure alone do not supply cancellation. The arithmetic inputs used in the component profile and in the following appendix must remain visible.
